\documentclass[11pt]{article}


\setlength{\parindent}{0pt}
\usepackage{xltxtra}
\usepackage{hyperref}
\hypersetup{hidelinks}
\usepackage{url}
\urlstyle{tt}
\usepackage{xcolor}
\definecolor{CVBlue}{RGB}{23,110,191}
\usepackage{calc}
\usepackage{graphicx}
\usepackage{tikz}
\usetikzlibrary{calc}
\usepackage{fontspec}
\usepackage{xeCJK}
\usepackage{enumitem}
\CJKsetecglue{} %% 取消中文与数字之间的间隙


%% 主文档字体设置
\setmainfont[
    Path = fonts/Main/,
    Extension = .otf,
    BoldFont = texgyretermes-bold.otf, % 加粗字体
]{texgyretermes-regular.otf} % 正文字体

% 中文字体设置
\setCJKmainfont[
    Path = fonts/hansans/,
    Extension = .ttf,
    BoldFont = NotoSansSC-Bold.ttf, % 加粗字体
]{NotoSansSC-Regular.otf} % 正文字体


\usepackage{relsize}
\usepackage{xspace}

% 使用 fontawesome（部分图标）
\usepackage{fontawesome} 

% A4纸，上下左右边距
\usepackage[
    a4paper,
    left=1.2cm,
    right=1.2cm,
    top=1.5cm,
    bottom=1cm,
    nohead
]{geometry}

\renewcommand{\baselinestretch}{1.5} % 行间距设为1.5

\usepackage{titlesec}
\usepackage{enumitem}
\setlist{noitemsep} % 取消列表项间的额外间距
%\setlist{nosep} % 取消所有垂直间距
\setlist[itemize]{topsep=0.25em, leftmargin=*}
\setlist[enumerate]{topsep=0.25em, leftmargin=*}

% --- 用于控制【不同项目之间】的垂直距离 ---
\newlength{\interProjectSpacing}
\setlength{\interProjectSpacing}{0.9em} % <--- 在此调整项目之间的距离
\newcommand{\projectsep}{\vspace{\interProjectSpacing}}

% --- 用于控制【项目标题】与下方【项目描述】的距离 ---
\newlength{\intraProjectTitleSep}
\setlength{\intraProjectTitleSep}{0.4em} % <--- 在此调整标题和描述的距离
\newcommand{\titlebreak}{\\[\intraProjectTitleSep]}

% --- 用于控制【项目描述】与下方【要点列表】的距离 ---
\newlength{\intraProjectListTopSep}
\setlength{\intraProjectListTopSep}{0.2em} % <--- 在此调整描述和列表的距离

% =======================================================================


\titleformat{\section}         % 定制 \section 命令 
{\large\bfseries\raggedright} % 将 section 标题设置为大号、粗体且左对齐
{}{0em}                      % 可用于为所有 section 添加前缀（如“章节...”）
{}                           % 可用于在标题前插入代码
[{\color{CVBlue}\titlerule}]  % 在标题后插入一条横线
\titlespacing*{\section}{0cm}{*1.6}{*1.2}



\begin{document}
\pagenumbering{gobble}

%%%% 利用tikz来定位照片
\begin{tikzpicture}[remember picture, overlay] 
    \node[anchor = north east] at ($(current page.north east)+(-2cm,-1.2cm)$) {\includegraphics[height=3cm]{avatar.jpg}};
  \end{tikzpicture}%
  %%%% 利用tikz来定位学校Logo，这里只在第一页显示
  \begin{tikzpicture}[remember picture, overlay] 
    \node[anchor = north west] at ($(current page.north west)+(0.5cm,+1.0cm)$) {\includegraphics[height=6cm]{zju.png}};
  \end{tikzpicture}%
\centerline{\LARGE\bfseries{崔雨纯}} 

\centerline{\normalsize{\faPhone\ 111-1111-1111 \quad \faEnvelopeO\ \href{mailto:3230100000@zju.edu.cn}{3230100000@zju.edu.cn}}} 

\centerline{\normalsize{\faGithubSquare\ \href{https://github.com/maksymilan}{https://github.com/maksymilan} \quad \faRssSquare\ \href{https://maksymilan.github.io/}{https://maksymilan.github.io}}} 
    
\section{\makebox[\widthof{\faGraduationCap}][c]{\color{CVBlue}\faGraduationCap}\ 教育背景}    
\textbf{中山大学} \hfill 2023.9 -- 至今\\[0.5em] % 标题和正文间加一点距离
光电信息科学与工程专业\quad 大三 
\begin{itemize}[nosep]
    \item 相关课程：《物理光学》、《半导体物理》、《显示科学与技术》
\end{itemize}

\section{\makebox[\widthof{\faUsers}][c]{\color{CVBlue}\faUsers}\ 项目经历}

% --- 第一个项目 ---
\textbf{TVPro 智能交互小电视} \hfill 2025.12 -- 至今 \titlebreak

项目描述：为实现边缘端智能交互，本项目基于 ESP32 平台，整合大模型与智能硬件，开发了一款可在线运行的 AI 交互终端，实现从语音输入到硬件响应的全链路闭环。

% 在 itemize 的选项中，使用 topsep=\intraProjectListTopSep 来控制上边距
\begin{itemize}[nosep, topsep=\intraProjectListTopSep]
    \item \textbf{开发环境与工具链}: 全程使用 \textbf{VSCode + PlatformIO} 完成项目开发、调试与固件烧录，实现了高效的嵌入式开发与部署流程。
    \item \textbf{大模型接入与交互}: 完成讯飞实时语音识别、星火大模型及 CosyVoice TTS 模型的全流程接入，通过 \textbf{WebSocket} 流式传输实现低延迟对话，将自然语言指令解析为可执行的硬件控制逻辑。
    \item \textbf{智能硬件驱动与交互}: 基于 \textbf{I2S} 协议实现 INMP441 全向麦克风音频采集与 Max98357 功放模块音频输出；通过 GPIO 驱动 ST7789 1.3 寸屏幕，实时显示时间、天气及交互状态，完成从“听”到“说”再到“显示”的硬件闭环。
    \item \textbf{网络通信与系统扩展}: 利用 \textbf{HTTPS/WebSocket} 完成云端 API 请求与设备状态同步；设计双固件（TV-PRO / 小智 AI）架构，支持一键切换，便于二次开发与功能扩展。
    
\end{itemize}

% 使用 \projectsep 命令来分隔两个项目
\projectsep

% --- 第二个项目 ---
\textbf{EDA-Camera 卡片相机} \hfill 2026.02 -- 至今 \titlebreak

项目描述：基于 ESP32S3 平台，开发了一款支持 500 万像素 CMOS 采集的卡片相机，通过多帧缓存与双画布缓冲技术，实现了高帧率实时预览与局域网相册管理。

% 在 itemize 的选项中，使用 topsep=\intraProjectListTopSep 来控制上边距
\begin{itemize}[nosep, topsep=\intraProjectListTopSep]
    \item \textbf{图像采集与处理}: 支持 OV2640/OV5640 模块化 CMOS 传感器，实现最高 2.5K 分辨率图像采集；通过 TJpg\_Decoder 库完成 JPEG 到 RGB565 格式转换，结合 7 种色彩滤镜及亮度/对比度/饱和度调节功能，完成端侧图像处理。
    \item \textbf{实时系统与显示优化}: 基于 \textbf{FreeRTOS Kernel V10.4.3} 实现多核多任务并发处理，封装 cameraTask 等任务库统一管理；采用 TFT\_eSPI 驱动 ST7789 2.4 寸屏幕，通过双画布缓冲机制避免画面重叠闪烁，保障预览流畅度。
    \item \textbf{存储与交互设计}: 实现局域网相册管理功能，支持无需读卡器直接导出图像至手机等设备；采用 14500 锂电池与 TP4056 充电方案，保障设备便携性与续航能力。
    \item \textbf{开发环境与工具链}: 全程使用 \textbf{VSCode + PlatformIO}，基于 C/C++ 语言完成项目开发、调试与固件烧录，实现了高效的嵌入式开发与部署流程。
\end{itemize}

% --- 第三个项目 ---
\textbf{基于多基色光源的显示-接收色域链路分析与图像优化研究} \hfill 202X.XX -- 202X.XX \titlebreak

项目描述：面向宽色域高保真显示技术研究需求，开展多基色显示系统的色域扩展与色彩传递链路优化研究，突破传统RGB三基色色域限制，建立显示-接收端色彩一致性模型，实现图像色彩的精准还原与优化。

\begin{itemize}[nosep, topsep=\intraProjectListTopSep]
    \item \textbf{理论建模与仿真分析}: 基于CIE 1931/2006色彩匹配函数，搭建多基色显示系统光谱色域分析模型，通过\textbf{MATLAB}完成数值仿真，量化分析第四基色波段对色域覆盖率、色域面积的提升效果，对比不同色彩匹配函数下的色域扩展差异。
    \item \textbf{色彩链路构建与算法开发}: 结合相机光谱响应模型，构建从显示端输出到接收端成像的完整\textbf{显示-接收色域传递链路}，编写MATLAB代码实现图像色彩重映射与优化算法，完成算法小程序的开发与调试。
    \item \textbf{实验验证与结果优化}: 基于实验室显示特性测量平台获取显示器光谱功率分布数据，依托色彩数据库/传感器参数库完成模型验证，分析仿真数据并总结多基色系统色域优化方法，为多基色显示光谱设计提供理论依据。
\end{itemize}

\section{\makebox[\widthof{\faCogs}][c]{\color{CVBlue}\faCogs}\ 技术栈}
\begin{itemize}[nosep]
    \item \textbf{编程语言：} \textbf{C/C++}, Python, Go, MATLAB
    \item \textbf{开发工具：} VSCode, PlatformIO, Git, SSH, LaTeX, MakeFile
    \item \textbf{操作系统：} Linux
\end{itemize}
\section{\makebox[\widthof{\faGraduationCap}][c]{\color{CVBlue}\faList}\ 获奖情况}
\begin{itemize}
    \item CC98年度用户 \hfill 2023.12
    
\end{itemize}
    
\section{\makebox[\widthof{\faInfo}][c]{\color{CVBlue}\faInfo}\ 其他}
\begin{itemize}[parsep=0.5ex]
    \item \textbf{GitHub：} \href{https://github.com/maksymilan}{https://github.com/cycist} 
    \item \textbf{英语水平：} CET-4, CET-6
\end{itemize}
\end{document}
